\documentclass[11pt,a4paper]{article}
\usepackage[margin=1in]{geometry}
\usepackage{hyperref}
\usepackage{enumitem}
\usepackage{graphicx}
\usepackage{xcolor}

% -----------------------------------------------------
% bvraghav-tiet-ucs503p-article.sty
% -----------------------------------------------------

% Variable: Lab Instructor
\makeatletter \newcommand{\BvrSetLabInstructor}[1]{%
  \gdef\Bvr@LabInstructor{#1}}
% default
\newcommand{\Bvr@LabInstructor}{Paramveer Sidhu}
% public accessor
\newcommand{\BvrLabInstructorValue}{\Bvr@LabInstructor}
\makeatother

% Add subtitle
\usepackage{titling}
% -- Set title bold
\pretitle{\begin{center}\bfseries\fontsize{18bp}{18bp}\selectfont}
\makeatletter
\newcommand{\BvrSubtitle}[1]{%
  \posttitle{%
    \par\end{center}
    \begin{center}\large#1\end{center}
    \vskip 0.5em}%
}
\makeatother

% Hints in the section title(s)
\newcommand{\BvrHint}[1]{%
  {\normalfont\normalsize\itshape\hfill #1}}

% -----------------------------------------------------

\title{PatchWatch: An Automated Dependency Vulnerability and Update Monitoring System}

\BvrSetLabInstructor{Paramveer Sidhu}

\BvrSubtitle{%
  Project Proposal for UCS503P  \\
  Submitted to: \BvrLabInstructorValue \\ \bfseries
  Thapar Institute of Engineering and Technology% 
}

\author{
  Diya Gupta, CSED%
  \thanks{Roll No: \texttt{1024030004}}
  \and 
  Mysha Bansal, CSED%
  \thanks{Roll No: \texttt{1024030010}}
  \and 
  Durva Garg, CSED%
  \thanks{Roll No: \texttt{1024030007}}
}

\date{\today}

\begin{document}
\maketitle

\section{Higher-order goal%
  \BvrHint{\ldots secondary framing}}
This project aims to improve software supply-chain
hygiene by reducing the time and effort developers spend
tracking, evaluating, and applying dependency updates.
While the underlying problem (dependency drift and
unpatched vulnerabilities) is broad, the immediate
engineering objective is to translate \emph{actionable
  dependency risk awareness} into a measurable, deployable
software solution.

\section{Time-to-value %
\BvrHint{\ldots primary heuristic}}
\begin{itemize}[leftmargin=0.7cm]
    \item We will deliver meaningful increments quickly by prototyping the core scan-and-alert workflow (scan a repository, detect outdated/vulnerable dependencies, surface them) and validating it against real open-source repositories within a short iteration window.
    \item We will prioritize fast feedback over perfection by using weekly iterations (and targeting CI-backed automation early) to reduce release friction.
    \item Success is defined by what can be delivered and measured in the first iteration -- a working scanner and dashboard -- then improved based on feedback from test repositories and users.
\end{itemize}

\section{Problem Statement}
Developers maintaining real-world projects struggle to keep dependencies both up to date and secure. Common pain points include:
\begin{itemize}[leftmargin=0.7cm]
    \item \textbf{Fragmentation:} dependency health information is scattered across registries (npm, PyPI, Maven), vulnerability databases (OSV, NVD, GitHub Advisory Database), and changelogs, with no single actionable view.
    \item \textbf{Low signal-to-noise:} generic ``newer version available'' tools flag every update equally, causing alert fatigue and making developers ignore genuinely critical security patches.
    \item \textbf{No project-specific impact analysis:} existing tools rarely tell a developer \emph{whether their own code} actually uses the affected API/behavior, or which files would need changes.
    \item \textbf{Slow remediation:} even when an update is known to be needed, manually editing manifests, verifying, and testing the change takes time that delays patching.
\end{itemize}

\textbf{Why this matters now (contextual relevance):} Supply-chain attacks and dependency-based vulnerabilities (e.g., prototype pollution, deserialization flaws) are an increasingly common attack vector; a lightweight tool that turns a firehose of registry/vulnerability data into a prioritized, project-aware action list can meaningfully reduce the window of exposure.

\section{Proposed Solution %
\BvrHint{\ldots Software Proposition}}
\subsection{Overview}
Build a \textbf{web-based} ``PatchWatch'' system --- a lightweight, Dependabot-style tool that watches a project's dependencies, detects new releases and security patches, checks whether they matter to the project, and tells the developer what to update and why. Core components:
\begin{itemize}[leftmargin=0.7cm]
    \item \textbf{Dependency scanner:} parses manifest files (\texttt{package.json}, \texttt{requirements.txt}, \texttt{pom.xml}) to extract current dependency versions.
    \item \textbf{Version monitor:} periodically checks package registries (npm, PyPI, Maven) for newer releases.
    \item \textbf{Security intelligence:} cross-references dependencies against vulnerability databases (OSV, NVD, GitHub Advisory Database).
    \item \textbf{Update classification:} labels each available update as patch/minor/major and flags security-relevant ones.
    \item \textbf{Impact analysis:} performs a lightweight static scan of the codebase to identify which files actually use the affected package/API, so alerts are project-specific rather than generic.
    \item \textbf{Dashboard and GitHub integration:} presents a prioritized view of dependency risk and can open a pull request for a recommended, low-risk update.
\end{itemize}

\subsection{Core workflow}
\begin{enumerate}[leftmargin=0.7cm]
    \item Developer connects a GitHub repository; PatchWatch scans manifest files and builds a dependency list.
    \item PatchWatch checks each dependency against the relevant package registry (new version?) and vulnerability databases (known CVE/advisory affecting the current version?).
    \item The update analyzer classifies each available update (patch/minor/major) and determines security severity.
    \item The impact analyzer greps/parses the codebase to find files that import or call the affected package, estimating project impact and required migration effort.
    \item PatchWatch produces a ranked dashboard entry and, optionally, opens an automated pull request with the recommended version bump and a summary of why the change is safe.
\end{enumerate}

\subsection{Operational constraints for an educational, capstone setting}
\begin{itemize}[leftmargin=0.7cm]
    \item Initially support a small set of ecosystems (JavaScript/npm, Python/pip, Java/Maven) rather than trying to cover every package manager.
    \item Avoid dependence on advanced research algorithms (no ML); focus on registry/vulnerability-DB integration, classification heuristics, and lightweight static analysis.
    \item Design for deployability using common web hosting, a managed database, and background job workers for periodic scans.
\end{itemize}

\section{Solution Approach %
\BvrHint{\ldots Engineering focus}}
This is an engineering course project, so we focus on scalable data-integration and workflow aspects rather than novel algorithm research.

\subsection{Update classification and impact analysis %
\BvrHint{\ldots non-research}}
Classification and impact detection will be heuristic-based, not ML-based:
\begin{itemize}[leftmargin=0.7cm]
    \item Use semantic-version parsing to classify an update as patch, minor, or major.
    \item Cross-reference the current version range against vulnerability database entries (OSV/NVD/GitHub Advisory) to flag security-relevant updates and assign a severity (low/medium/high/critical).
    \item Use simple static analysis (import/require statements, Tree-sitter or language ASTs) to list files referencing the affected package, without claiming complete or state-of-the-art code understanding.
    \item Always present recommendations to the developer for review; never auto-merge without explicit approval.
\end{itemize}

\subsection{Web architecture %
\BvrHint{\ldots web-based solution}}
A typical 3-tier web architecture with an added background-processing layer:
\begin{itemize}[leftmargin=0.7cm]
    \item \textbf{Frontend:} responsive dashboard for viewing dependency risk, drilling into a specific package, and triggering PR creation.
    \item \textbf{Backend API:} authentication (GitHub OAuth), repository management, scan orchestration, classification/impact endpoints.
    \item \textbf{Background jobs:} scheduled workers that poll package registries and vulnerability databases and re-scan repositories.
    \item \textbf{Database:} stores repositories, dependency snapshots, scan history, vulnerability matches, and PR status.
\end{itemize}

\subsection{CI/CD and fast delivery enabler}
CI/CD will be used to:
\begin{itemize}[leftmargin=0.7cm]
    \item Automatically build and run tests on every change to the PatchWatch codebase itself.
    \item Produce repeatable deployment artifacts to staging.
    \item Enable safe and frequent releases of incremental features (new ecosystem support, new vulnerability sources, etc.).
\end{itemize}

\section{Evaluation Criterion %
\BvrHint{\ldots measurable and attributable}}
We define success using measurable metrics tied to the core scan-detect-alert workflow.

\subsection{Primary evaluation metric}
\textbf{Time-to-detection (TTD):} time between a security advisory/patch being published for a dependency and PatchWatch surfacing an alert for a repository that uses the affected version.
\begin{itemize}[leftmargin=0.7cm]
    \item Target: median TTD $\le$ 24 hours for repositories under active scanning during the pilot window.
    \item Attribution: measured from advisory publish timestamp (per OSV/NVD/GitHub Advisory) vs. PatchWatch alert-creation timestamp.
\end{itemize}

\subsection{Secondary evaluation metrics}
\begin{itemize}[leftmargin=0.7cm]
    \item \textbf{Alert precision:} percentage of security alerts that are relevant (i.e., the flagged package/API is actually used in the codebase), reducing noise compared to naive ``new version available'' tools.
    \item \textbf{PR acceptance rate:} percentage of automatically generated update pull requests that pass CI and are merged without manual edits, on pilot test repositories.
    \item \textbf{Coverage:} number/percentage of dependencies across supported ecosystems successfully scanned and classified per repository.
    \item \textbf{Reliability:} availability of the scanning service and dashboard during business hours (e.g., $\ge$ 99\% uptime in pilot).
\end{itemize}

\subsection{Pilot validation plan %
\BvrHint{\ldots high-level}}
\begin{itemize}[leftmargin=0.7cm]
    \item Connect PatchWatch to 10--15 pilot repositories (a mix of team projects and public open-source repositories with known historical vulnerabilities).
    \item Compare PatchWatch's detection timing and alert quality against manually checking registries/advisories for the same repositories over the iteration window.
\end{itemize}

\section{Scalability %
\BvrHint{\ldots with foresight}}
The proposition should scale in both theoretical and practical senses.

\begin{enumerate}[leftmargin=0.7cm]
    \item \textbf{Scalable (theoretically):} The system should support growth in number of monitored repositories and dependencies by:
    \begin{itemize}[leftmargin=0.7cm]
        \item decoupling scanning/analysis workers from the API and dashboard,
        \item queuing and rate-limiting registry/vulnerability-DB requests (e.g., via Redis-backed job queues),
        \item enabling pagination and indexing for common dashboard queries.
    \end{itemize}
    
    \item \textbf{Operationally deployable (practically):} The system should run on common web hosting:
    \begin{itemize}[leftmargin=0.7cm]
        \item managed relational database,
        \item containerized or platform-hosted deployment,
        \item CI/CD pipeline for staging/production.
    \end{itemize}
\end{enumerate}

\section{Engine availability heuristic}
A mature set of ``engines'' (tooling/services) is available to implement this as a web-based project:
\begin{itemize}[leftmargin=0.7cm]
    \item Web framework for REST APIs and a reactive dashboard (Node.js/Express backend, React/TypeScript frontend).
    \item Managed relational database (PostgreSQL) with an ORM (Prisma).
    \item Background job queue for scheduled scans (Redis + BullMQ).
    \item Public registries and databases: npm/PyPI/Maven registry APIs, OSV API, NVD, GitHub Advisory Database.
    \item GitHub REST API and OAuth for repository access and PR creation.
    \item Lightweight code-analysis tooling (Tree-sitter / language ASTs) for impact detection.
    \item Test framework for unit/integration tests; CI runner and staging deployment target (e.g., Vercel + Render/Railway, Docker).
\end{itemize}
We will use these mature components to focus on data integration, classification logic, and measurement rather than inventing new infrastructure.

\section{Project scope and deliverables %
\BvrHint{\ldots iteration-friendly}}
\subsection{Initial deliverable %
\BvrHint{\ldots first iteration}}
\begin{itemize}[leftmargin=0.7cm]
    \item Dependency scanner for one ecosystem (npm/\texttt{package.json}) extracting current versions.
    \item Version monitor querying the npm registry for latest versions.
    \item Basic dashboard listing dependencies with current vs. latest version and simple patch/minor/major classification.
    \item Basic logging and timestamps for TTD measurement.
    \item CI: build + backend unit tests + linting.
    \item CD to staging with a single command or pipeline job.
\end{itemize}

\subsection{Subsequent deliverables}
\begin{itemize}[leftmargin=0.7cm]
    \item Security intelligence integration (OSV/NVD/GitHub Advisory Database) with severity scoring.
    \item Impact analysis: static-analysis pass identifying files/usages affected by a given dependency update.
    \item GitHub OAuth integration, repository scanning, and automatic pull-request creation for recommended updates.
    \item Support for additional ecosystems (Python/\texttt{requirements.txt}, Java/\texttt{pom.xml}).
    \item Notification layer and richer admin dashboard (backlog of unresolved security issues, historical trends).
\end{itemize}

\section{Risks and mitigations}
\begin{itemize}[leftmargin=0.7cm]
    \item \textbf{Third-party API rate limits:} registry and vulnerability-DB APIs may throttle requests; mitigate with caching, batching, and a background job queue rather than synchronous on-demand calls.
    \item \textbf{False positives/negatives in impact analysis:} static analysis may miss dynamic usages or over-flag files; mitigate by presenting impact analysis as an aid for developer review, not an automatic gate.
    \item \textbf{Data staleness:} vulnerability databases and registries update independently; mitigate with a defined re-scan interval and clear ``last checked'' timestamps.
    \item \textbf{Team adoption of auto-generated PRs:} developers may distrust automated updates; mitigate by keeping PRs low-risk by default (patch/security-only), always running tests, and never auto-merging.
\end{itemize}

\section{Summary}
This project proposes PatchWatch, a deployable web platform that reduces the time and effort required to detect, prioritize, and remediate outdated or vulnerable project dependencies. It meets the course criteria by emphasizing time-to-value through rapid iterations, implementing CI/CD to support frequent safe releases, and using measurable evaluation criteria (especially time-to-detection and alert precision). It remains focused on engineering integration, classification, and workflow design rather than research-driven novelty, while still offering a substantial, demonstrable end-to-end system (scanner $\rightarrow$ security intelligence $\rightarrow$ impact analysis $\rightarrow$ dashboard $\rightarrow$ automated PR).

\end{document}
